% ==================== 6 结论 ====================
\section{结　论}

本文围绕 8 组 200--900 $^\circ$C 热处理的 TiO$_2$ 样品，完成了作业二要求的两件事。数据处理方面，先用 \texttt{make\_\allowbreak origin\_\allowbreak ready.py} 把仪器导出的 9 行表头 txt 改成 Origin 能稳定识别的 3 行表头 + Tab 分隔格式，解决了逗号分隔符和表头误判两个问题；再按样品各自归一化到最强峰，在 Origin 中用 Stacked Lines by Y Offsets 画出 8 条曲线，并把锐钛矿 71-1167 和金红石 75-1756 两张 PDF 卡片用 Stackplot + 新 Layer 叠加到主图下方。最终的图能同时看到 8 组样品的峰位变化和两套标准卡片的对应关系。

物相方面，200 $^\circ$C 的谱线只有一个宽化的锐钛矿 (101) 包络，结晶度很低；300--600 $^\circ$C 都能用锐钛矿卡片指标化，25.4$^\circ$ (101)、37.8$^\circ$ (004)、48.0$^\circ$ (200)、54.0$^\circ$ (105)、55.1$^\circ$ (211) 是主要特征峰；700 $^\circ$C 在 27.6$^\circ$ 出现金红石 (110)，说明相变开始；800 $^\circ$C 两相共存，金红石的 (110)、(101)、(211) 都很明显；900 $^\circ$C 只剩金红石。这个趋势和 TiO$_2$ 热处理过程中锐钛矿向金红石转变的常识一致。本文只做定性描述，未做更深入的定量分析，这部分留到后续课程作业。

实验操作方面，制样和光路对中是数据质量的第一道关：样品要磨细、量足、压平、与槽沿齐平，每次换样都要“校读”；开机先水后高压、关机先降压再关束，升降压要慢；测量过程中不能开门，角度范围不要随意超出 20$^\circ$--80$^\circ$。数据处理方面，原始 txt 不要直接拖进 Origin，先用脚本转成 3 行表头 + Tab；PDF 卡片要核对三根最强线，不能只看卡片编号；卡片层的 X 轴必须和主图一致，位置和大小要按页面比例调。本次实验遇到的低角背底高、峰位偏移、两相峰重叠、Origin 表头错位、Jade 卡片选错等问题，以及对应的解决办法，都整理在第 5 节的表~\ref{tab:problems} 中。

通过这次作业，我对“原始数据—格式转换—多曲线绘图—PDF 卡片对照—物相判断”这条流程有了具体的体会，也知道了哪些步骤可以脚本化、哪些步骤必须回到实验现场去判断。后续如果继续做这个体系，可以补测标样、延长扫描时间、把 700--800 $^\circ$C 细分成更多温度点，并尝试峰拟合和定量分析；本次作业先把数据整理、绘图和实验总结做扎实。
